% Phụ lục D — Điều khoản dịch vụ. Viết để đọc độc lập: không dùng mã vai trò
% hệ thống, mã yêu cầu chức năng hay tham chiếu chéo vào thân báo cáo.

\section*{Phụ lục D. Điều khoản dịch vụ}
\addcontentsline{toc}{section}{Phụ lục D. Điều khoản dịch vụ}

\noindent\textbf{Điều 1. Phạm vi áp dụng}\\
Điều khoản này áp dụng cho mọi tổ chức, cá nhân sử dụng nền tảng Farmery,
bao gồm nhà cung cấp nông sản, doanh nghiệp mua hàng với số lượng lớn và
người tiêu dùng mua hàng cho nhu cầu cá nhân. Việc đăng ký tài khoản và sử
dụng nền tảng đồng nghĩa với việc chấp nhận toàn bộ các điều khoản dưới đây.

\noindent\textbf{Điều 2. Tài khoản}\\
Mỗi tài khoản chỉ được sử dụng đúng với mục đích đã đăng ký và phạm vi quyền
hạn được Farmery cấp. Người dùng chịu trách nhiệm bảo mật thông tin đăng
nhập và chịu trách nhiệm về toàn bộ hoạt động phát sinh từ tài khoản của
mình; khi nghi ngờ tài khoản bị người khác truy cập trái phép, người dùng
phải thông báo ngay cho Farmery.

\noindent\textbf{Điều 3. Quyền và nghĩa vụ của nhà cung cấp}\\
Nhà cung cấp cam kết cung cấp thông tin sản phẩm, giấy chứng nhận và nhật ký
canh tác trung thực, chính xác; chịu trách nhiệm về chất lượng hàng hóa đúng
như đã mô tả; cập nhật kịp thời số lượng hàng còn lại và giá bán; giao hàng
đúng cam kết về thời gian và điều kiện bảo quản đối với hàng dễ hư hỏng.
Nhà cung cấp có quyền phản hồi công khai đối với đánh giá của người mua.

\noindent\textbf{Điều 4. Quy định về hàng hóa được phép đăng bán}\\
Farmery là nền tảng chuyên biệt cho nông sản. Nhà cung cấp chỉ được đăng bán
hàng hóa thuộc các nhóm sau: nông sản tươi chưa qua chế biến; nông sản đã sơ
chế hoặc bảo quản đơn giản mà không thêm phụ gia; và nông sản chế biến sâu,
với điều kiện bắt buộc phải có bản tự công bố sản phẩm cùng giấy chứng nhận
cơ sở đủ điều kiện an toàn thực phẩm còn hiệu lực.

Không được đăng bán trên nền tảng: thủy sản sống, thịt tươi và động vật
sống; thực phẩm chức năng; thuốc bảo vệ thực vật và vật tư nông nghiệp; mọi
hàng hóa phi thực phẩm; và mọi hàng hóa bị cấm kinh doanh theo quy định của
pháp luật.

Mỗi sản phẩm phải được gắn đúng vào danh mục do Farmery quy định, kèm đầy đủ
thông tin bắt buộc của danh mục đó. Nhà cung cấp đăng sản phẩm sai danh mục
hoặc thuộc nhóm không được phép sẽ bị xử lý theo Điều 8.

\noindent\textbf{Điều 5. Quyền và nghĩa vụ của người mua}\\
Người mua có quyền tra cứu thông tin nguồn gốc hàng hóa trước khi giao dịch,
được bảo vệ theo chính sách đổi trả và hoàn tiền tại Điều 7, và được gửi
đánh giá sau khi đã xác nhận nhận hàng. Người mua có nghĩa vụ thanh toán
đúng hạn, cung cấp thông tin nhận hàng chính xác, và phản hồi trung thực khi
gửi khiếu nại hoặc đánh giá.

\noindent\textbf{Điều 6. Phân định trách nhiệm theo hình thức bán hàng}\\
Nền tảng vận hành hai hình thức bán hàng với trách nhiệm khác nhau:
\begin{itemize}
    \item \textbf{Hàng do nhà cung cấp trực tiếp bán:} nhà cung cấp là bên
    bán trong giao dịch và chịu trách nhiệm về chất lượng, số lượng cũng như
    tiến độ giao hàng. Farmery giữ vai trò trung gian kết nối, bảo đảm minh
    bạch thông tin và hỗ trợ giải quyết tranh chấp.
    \item \textbf{Hàng do Farmery mua lại rồi bán:} Farmery là bên bán trong
    giao dịch và trực tiếp chịu trách nhiệm với người mua về chất lượng hàng
    hóa cũng như việc đổi trả, không phụ thuộc vào trách nhiệm của nhà cung
    cấp ban đầu.
\end{itemize}
Thông tin về bên bán được hiển thị rõ trên trang sản phẩm và trên đơn hàng.

\noindent\textbf{Điều 7. Phí dịch vụ}\\
% TODO: nhóm điền mức phí cụ thể theo phân tích chi phí vận hành thực tế.
Đối với hàng do nhà cung cấp trực tiếp bán, Farmery thu phí hoa hồng trên
mỗi giao dịch thành công. Đối với hàng do Farmery mua lại rồi bán, Farmery
không thu phí của nhà cung cấp mà hưởng chênh lệch giữa giá mua và giá bán.
Ngoài ra, nhà cung cấp có thể tự nguyện sử dụng các dịch vụ trả phí gồm xác
thực chứng nhận nhanh và gói hiển thị nổi bật; vị trí hiển thị mua được luôn
gắn nhãn ``Tài trợ'' để phân biệt với kết quả tìm kiếm thông thường.

Farmery không thu phí đăng ký tài khoản và không thu phí duy trì tài khoản
đối với nhà cung cấp. Việc được ưu tiên hiển thị nhờ điểm uy tín cao là
quyền lợi miễn phí, hoàn toàn độc lập với các dịch vụ trả phí nói trên.

\noindent\textbf{Điều 8. Thanh toán}\\
Giao dịch với người tiêu dùng được thanh toán qua ví điện tử, chuyển khoản
ngân hàng hoặc trả tiền khi nhận hàng.

Đối với giao dịch số lượng lớn có giá trị vượt ngưỡng do Farmery công bố,
tiền của bên mua được một đơn vị trung gian thanh toán được cấp phép giữ
lại. Farmery không trực tiếp giữ tiền của người dùng. Khoản tiền này chỉ
được chuyển cho bên bán sau khi bên mua xác nhận đã nhận đủ hàng đúng cam
kết. Khi phát sinh tranh chấp, khoản tiền được giữ nguyên tại đơn vị trung
gian cho tới khi khiếu nại được giải quyết xong theo Điều 9.

\noindent\textbf{Điều 9. Đổi trả và hoàn tiền}\\
% TODO: nhóm điền khung thời gian gửi khiếu nại và thời gian xử lý tối đa
% sau khi chốt được năng lực vận hành thực tế.
Tùy mức độ và nguyên nhân xác định được, khiếu nại được giải quyết theo một
trong các hướng: hoàn tiền toàn phần, hoàn tiền một phần, đổi sang lô hàng
khác, hoặc từ chối khiếu nại nếu bằng chứng không đầy đủ. Do đặc thù hàng
nông sản dễ hư hỏng, người mua cần gửi khiếu nại kèm hình ảnh minh chứng
trong thời hạn Farmery công bố kể từ khi nhận hàng.

\noindent\textbf{Điều 10. Kiểm duyệt và xử lý vi phạm}\\
% TODO: nhóm điền thời gian xử lý hồ sơ xác minh doanh nghiệp sau khi chốt
% quy trình thẩm định thực tế.
Farmery kiểm duyệt hồ sơ nhà cung cấp, sản phẩm và giấy chứng nhận trước khi
cho hiển thị công khai, đồng thời kiểm tra lại định kỳ trong quá trình hoạt
động. Vi phạm được xử lý theo mức độ tăng dần: nhắc nhở, ẩn sản phẩm, cảnh
cáo kèm trừ điểm uy tín, tạm khóa và cuối cùng là khóa vĩnh viễn tài khoản.

Các hành vi bị xử lý khóa tài khoản ngay lập tức, không qua bước trung gian,
gồm: giả mạo giấy chứng nhận, gian lận trong giao dịch, và thao túng hệ
thống đánh giá.

\noindent\textbf{Điều 11. Sở hữu trí tuệ và dữ liệu}\\
Dữ liệu nhật ký canh tác, hình ảnh sản phẩm và giấy chứng nhận do nhà cung
cấp đưa lên vẫn thuộc quyền sở hữu của nhà cung cấp. Farmery được quyền sử
dụng dữ liệu này để vận hành nền tảng, hiển thị công khai trang thông tin
nguồn gốc sản phẩm, và cải thiện các tính năng hỗ trợ người dùng; với mục
đích cuối, dữ liệu có thể được sử dụng ở dạng ẩn danh hoặc tổng hợp khi
không cần gắn với danh tính cụ thể.

\noindent\textbf{Điều 12. Giới hạn trách nhiệm}\\
Đối với hàng do nhà cung cấp trực tiếp bán, trách nhiệm của Farmery giới hạn
trong phạm vi thông tin đã được xác minh qua quy trình kiểm duyệt; các tranh
chấp về chất lượng nằm ngoài phạm vi xác minh được giải quyết theo Điều 9.
Farmery không tự sản xuất và không tự vận chuyển hàng hóa; việc vận chuyển
do các đơn vị dịch vụ độc lập thực hiện.

\noindent\textbf{Điều 13. Chấm dứt hoặc tạm ngừng dịch vụ}\\
Người dùng có quyền tự nguyện đóng tài khoản bất kỳ lúc nào, với điều kiện
không còn đơn hàng hoặc hợp đồng đang thực hiện dở dang. Farmery có quyền
tạm ngừng hoặc chấm dứt tài khoản theo Điều 10, hoặc khi phát hiện dấu hiệu
gian lận nghiêm trọng. Việc chấm dứt không làm mất đi các nghĩa vụ tài chính
đã phát sinh trước đó, bao gồm các khoản đang được giữ tại đơn vị trung gian
thanh toán theo Điều 8.

\noindent\textbf{Điều 14. Giải quyết tranh chấp}\\
Tranh chấp giữa các bên trước hết được giải quyết qua thương lượng trực tiếp
trên kênh liên lạc của nền tảng. Nếu không đạt được thỏa thuận, các bên có
thể yêu cầu Farmery hòa giải dựa trên dữ liệu giao dịch và thông tin nguồn
gốc đã được hệ thống ghi nhận làm căn cứ khách quan. Nếu vẫn không giải
quyết được, tranh chấp được xử lý theo quy định của pháp luật Việt Nam.

\noindent\textbf{Điều 15. Thay đổi điều khoản}\\
Farmery có quyền cập nhật điều khoản dịch vụ khi cần, có thông báo trước cho
người dùng qua nền tảng hoặc qua địa chỉ liên lạc đã đăng ký. Việc tiếp tục
sử dụng nền tảng sau thời điểm điều khoản mới có hiệu lực được xem là chấp
nhận thay đổi.
